\section{Background and Related Work}

% IEEE target for 7-page paper: ~1.0 pages.
% This section should summarize prior work and establish what is new in this paper.
% Suggested structure:
% - defensive malware analysis methods
% - AI-assisted cybersecurity workflows
% - reproducibility and measurement gaps

% TODO: add 4--8 concise, citation-backed paragraphs.

Defensive malware analysis is typically evaluated across static and dynamic evidence streams.
Static analysis captures persisted artifacts such as metadata, strings, and code structure, while
dynamic analysis captures runtime behavior in isolated environments.

Recent AI-assisted security research has emphasized workflow support, code reasoning,
and analyst productivity. At the same time, less work exists on how controlled AI-assisted
conditions change the measurability of defensive outcomes.

This work differs in scope by avoiding claims about provenance and by treating workflow
assistance as an experimental condition. The emphasis is on defensible comparison, not on
building or improving offensive capability.
